\chapter{Uninformed learning with Tsallis-INF (Section 4.1)}
\label{chap:tsallis_inf}

\begin{lemma}[Runs of an uninformed learner]
  \label{lem:uninformed_run}
  \lean{Learning.IsAlgEnvSeq.lintegral_gameRound,
    Learning.IsAlgEnvSeq.lintegral_gameRound_of_weighted, Learning.IsAlgEnvSeq.ae_gameReward_mem,
    Learning.IsAlgEnvSeq.lintegral_gameAction_eq_sum,
    Learning.IsAlgEnvSeq.lintegral_one_sub_gameReward_mul,
    Ito2026Adversarial.policy_ofUninformed_tsallisINFHalf,
    Ito2026Adversarial.integral_one_sub_reward_run,
    Ito2026Adversarial.integral_importanceWeight_run,
    Ito2026Adversarial.integral_inner_halfEstimate_run, Ito2026Adversarial.integral_comp_action_run,
    Ito2026Adversarial.integral_stabWeight_run}
  \leanok
  \uses{def:repeated_game, def:tsallis_inf}
  In a run of Tsallis-INF against an adversary, $x_t$ has conditional law $p_t$ given the past
  rounds, $y_t$ and $x_t$ are conditionally independent given the past, and the estimate $g_t(x)$
  has conditional expectation $u(x, y_t)$ given the past and $y_t$. In particular
  $\E\langle e_x - p_t, g_t\rangle = \E[u(x, y_t) - u(x_t, y_t)]$ and $\E[f(x_t)] = \E[\sum_x
  p_t(x) f(x)]$ for $f \ge 0$.
\end{lemma}

\begin{proof}
  \leanok
  The policy of Tsallis-INF draws $x_t$ from $p_t$, a function of the past rewards; the adversary
  draws $y_t$ from its policy on the past rounds (simultaneous moves). Hence the integral of a
  nonnegative function of the history, $x_t$, $y_t$ and $r_t$ is the integral of $\sum_x p_t(x)
  \int\int F\, dR\, d\mathrm{opp}$. The reward has conditional mean $u(x_t, y_t)$, so $\E[g_t(x)
  \mid \text{past}, y_t] = 1 - p_t(x)(1 - u(x, y_t))/p_t(x)$; the integrability of the
  importance weights follows from these identities for nonnegative functions.
\end{proof}

\begin{lemma}[Strict PSNE: the self-bounding inequality]
  \label{lem:tsallis_strict_step}
  \lean{Ito2026Adversarial.externalRegretAgainst_le_of_isStrictPSNE,
    Ito2026Adversarial.sum_inv_sqrt_mul_sum_sqrt_le,
    Ito2026Adversarial.integral_sqrt_le_sqrt_integral}
  \leanok
  \uses{lem:tsallis_inf_self_bound, lem:regret_relations, lem:uninformed_run}
  In the setting of Theorem~\ref{thm:tsallis}, if $(x^\star, y^\star)$ is a strict PSNE, with $H_T
  = \sum_{t < T}\frac{1}{t + 1}$, $S = \sum_{x \ne x^\star}\frac{1}{\DR_x}$ and $C$ the expected
  number of rounds with $y_t \ne y^\star$,
  $\ER_T(x^\star) \le 12\sqrt{H_T S}\sqrt{\ER_T(x^\star) + 4C} + 6$.
\end{lemma}

\begin{proof}
  \leanok
  \uses{lem:tsallis_inf_self_bound, lem:regret_relations, lem:uninformed_run}
  By the Cauchy--Schwarz inequality, $\sum_t\frac{1}{\sqrt{t + 1}}\sum_{x \ne x^\star}
  \sqrt{p_t(x)} \le \sqrt{H_T S}\sqrt{\sum_t\langle\DR, p_t\rangle}$; by Jensen's inequality and
  Lemma~\ref{lem:uninformed_run}, $\E[\sqrt{\sum_t\langle\DR, p_t\rangle}] \le
  \sqrt{\E[\sum_t\DR_{x_t}]} \le \sqrt{\ER_T(x^\star) + 4C}$
  (Lemma~\ref{lem:regret_relations}~(3)); conclude with
  Lemma~\ref{lem:tsallis_inf_self_bound}.
\end{proof}

\begin{theorem}[Theorem 1]
  \label{thm:tsallis}
  \lean{Ito2026Adversarial.exists_psmr_tsallisINFHalf_le}
  \leanok
  \uses{def:tsallis_inf, def:regrets, lem:tsallis_inf_worst, lem:tsallis_inf_self_bound,
    lem:regret_relations, lem:sqrt_func, lem:self_bound, lem:uninformed_run,
    lem:tsallis_strict_step}
  There is a universal constant $C$ such that for every game $u$ with utilities in $[-1, 1]$, every
  adaptive adversary and every reward noise with conditional mean $u$ and rewards in $[-1, 1]$,
  Tsallis-INF with $\alpha = 1/2$ and $\eta_t = \frac{1}{2\sqrt t}$ has
  $\PSMR_T \le C \sqrt{m_x T}$; if $(x^\star, y^\star)$ is a strict PSNE,
  $\PSMR_T \le C (1 + \frac{1}{\DCmin}) \sum_{x \ne x^\star} \frac{1 + \log T}{\DR_x}$; and if
  $\Dmix > 0$, $\PSMR_T \le C \frac{m_x}{\Dmix}$.
\end{theorem}

\begin{proof}
  \leanok
  \uses{lem:tsallis_inf_worst, lem:tsallis_inf_self_bound, lem:regret_relations, lem:sqrt_func,
    lem:self_bound, lem:uninformed_run, lem:tsallis_strict_step}
  The Lean proof takes $C = 200$ (\texttt{psmr\_tsallisINFHalf\_le\_sqrt},
  \texttt{psmr\_tsallisINFHalf\_le\_of\_isStrictPSNE},
  \texttt{psmr\_tsallisINFHalf\_le\_of\_mixGap\_pos}).
  Worst case: $\PSMR_T \le \NR_T \le \ER_T \le 24\sqrt{m_x T} + 6$
  (Lemmas~\ref{lem:regret_relations} and~\ref{lem:tsallis_inf_worst}), and $6 \le 6\sqrt{m_x T}$ for
  $T \ge 1$. If $\Dmix > 0$: $\PSMR_T = \NR_T - \Dmix T \le 6 + 24\sqrt{m_x T} - \Dmix T \le 6 +
  \frac{144 m_x}{\Dmix}$ (Lemma~\ref{lem:sqrt_func}), and $6 \le \frac{12 m_x}{\Dmix}$ as $\Dmix \le
  2$. Strict PSNE: if $m_x = 1$, $\ER_T(x^\star) = 0$ and $\PSMR_T \le 0$. Otherwise
  Lemma~\ref{lem:tsallis_strict_step} and Lemma~\ref{lem:self_bound} give $\ER_T(x^\star) \le 144
  H_T S + 24\sqrt{H_T S C} + 12$ (if $\ER_T(x^\star) \ge 0$), and Lemma~\ref{lem:regret_relations}
  (2) and Lemma~\ref{lem:sqrt_func} in $C$ give $\PSMR_T \le \ER_T(x^\star) - \DCmin C \le 144 H_T S
  (1 + \frac{1}{\DCmin}) + 12$ (if $\DCmin = 0$, the adversary has a single action and $C = 0$).
  Finally $H_T \le 1 + \log T$ and $S \ge \frac12$ absorb the constant $12$.
\end{proof}

\begin{remark}
  \label{rem:tsallis_statement}
  The paper states the strict-PSNE bound as $O(\frac{1}{\DCmin}\sum_{x \ne x^\star}
  \frac{\log T}{\DR_x})$, which follows from the form above when $m_y \ge 2$ and $T \ge 2$ (then
  $\DCmin \le 2$). For $m_y = 1$ it is false in Lean, where $\DCmin$ is an infimum over no action,
  equal to $0$, and $C/0 = 0$. The paper states the third bound for games without PSNE, which need
  not have $\Dmix > 0$ (Remark~\ref{rem:no_psne_gap}).
\end{remark}
